\section*{अध्याय \ref{ch:b3:nervous-systems} --- तंत्रिका-तंत्रों का संगठन}

\begin{solution}{exo:b3:nervous-systems:1}
\omterm{def:b3:nervous-systems:cells}{तारा-कोशिकाएँ}: पोटैशियम और ग्लूटामेट के लिए बफ़र बनती हैं, न्यूरॉनों को
पोसती हैं, रक्त--मस्तिष्क अवरोध प्रेरित करती हैं, स्थानीय रक्त-प्रवाह
नियंत्रित करती हैं। \omterm{def:b3:nervous-systems:cells}{ऑलिगोडेंड्रोसाइट}: केंद्रीय तंत्रिका-तंत्र में
तंत्रिकाक्षों को \omterm{def:b3:nervous-systems:cells}{मायलिन} देती हैं। श्वान कोशिकाएँ: परिधीय तंत्रिकाक्षों
को \omterm{def:b3:nervous-systems:cells}{मायलिन} देती हैं (और उनके पुनर्जनन का मार्गदर्शन करती हैं)।
\omterm{def:b3:nervous-systems:cells}{सूक्ष्म-ग्लिया}: मस्तिष्क के \omterm{def:b3:innate-immunity:innate}{महाभक्षक}, जो तंत्रिका-संधियाँ छाँटती और
मलबा साफ़ करती हैं। (एपेंडिमल कोशिकाएँ निलयों की परत बनाती हैं और
\omterm{prop:b3:nervous-systems:barrier-budget}{प्रमस्तिष्क-मेरु द्रव} बनाती हैं।)
\end{solution}

\begin{solution}{exo:b3:nervous-systems:2}
उस रंजन ने कुछ न्यूरॉन पूरे-के-पूरे काले कर दिए और बाकी अदृश्य छोड़
दिए, जिससे कोई एक कोशिका पूरी दिख सकी। काहाल ने निष्कर्ष निकाला कि
न्यूरॉन अलग-अलग कोशिकाएँ हैं, जो बिना संलयित हुए एक-दूसरे को छूती हैं,
और जिनमें बहाव की दिशा द्रुमिकाओं से तंत्रिकाक्ष तक नियत है; गॉल्जी ने
निष्कर्ष निकाला कि वे कोई अविच्छिन्न जाल बनाती हैं। काहाल सही था;
इलेक्ट्रॉन सूक्ष्मदर्शी ने 1954 में वह संधि-अंतराल दिखा दिया।
\end{solution}

\begin{solution}{exo:b3:nervous-systems:3}
कोई ठोकर चतुःशिरस्क को खींचती है; उसके पेशी-तर्कु दगते हैं; Ia अभिवाही
\omterm{def:b3:nervous-systems:plan}{मेरुरज्जु} में घुसकर सीधे चतुःशिरस्क के प्रेरक न्यूरॉनों पर संधि बनाते
हैं, जो दगकर उस पेशी को सिकोड़ देते हैं, जबकि कोई \omterm{def:b3:nervous-systems:cells}{अंतर्न्यूरॉन} पश्चजंघ
के प्रेरक न्यूरॉनों का निरोध करता है। एक संधि, \qty{80}{m/s} पर छोटे
तंत्रिकाक्ष: लगभग \qty{30}{ms}। कोई ऐच्छिक लात मस्तिष्क से गुज़रती है
--- प्रत्यक्षण, निर्णय, प्रेरक नियोजन, अवरोही पथ, कई संधियाँ --- और
उसमें \qtyrange{200}{300}{ms} लगते हैं।
\end{solution}

\begin{solution}{exo:b3:nervous-systems:4}
अनुकंपी (लक्ष्य पर नॉरएड्रिनैलीन): हृदय तेज़ और प्रबल, पुतली फैली, आँत
की गतिशीलता तथा स्राव घटे, वायुमार्ग चौड़े। परानुकंपी (ऐसीटिलकोलीन):
हृदय धीमा, पुतली सिकुड़ी, आँत की गतिशीलता तथा स्राव बढ़े, वायुमार्ग
सँकरे।
\end{solution}

\begin{solution}{exo:b3:nervous-systems:5}
$\lambda = \sqrt{R_{m}d/4R_{i}} = \sqrt{10^{4}\times 10^{-4}/400}$ cm
$= \qty{0.05}{cm} = \qty{0.5}{mm}$। \qty{2}{mm} के बाद: $e^{-4} \approx
2\,\%$। $\lambda$ दुगुना करने के लिए चार गुना व्यास चाहिए, यानी
\qty{4}{\micro m}।
\end{solution}

\begin{solution}{exo:b3:nervous-systems:6}
संवेदी: $12\times 6 = \qty{72}{m/s}$, $1.2/72 = \qty{17}{ms}$। प्रेरक:
\qty{90}{m/s}, $1.1/90 = \qty{12}{ms}$। दो संधियाँ, \qty{1}{ms}: लगभग
\qty{30}{ms}। अमायलिनी: $1.1\sqrt{12} = \qty{3.8}{m/s}$ और
$1.1\sqrt{15} = \qty{4.3}{m/s}$: $315 + 258 + 1 \approx \qty{570}{ms}$
--- पैर बचाने के लिए बहुत धीमा।
\end{solution}

\begin{solution}{exo:b3:nervous-systems:7}
$2.4\times 10^{20}/8.6\times 10^{10} \approx 2.8\times 10^{9}$ ATP प्रति
न्यूरॉन प्रति सेकंड; प्रति आवेग $2\times 10^{9}$ पर, औसतन लगभग
\qty{1.4}{Hz}। कूट को विरल होना ही पड़ता है: थोड़े आवेग, और सूचना इसमें
कि कौन-से न्यूरॉन दगते हैं और कब, न कि हर जगह ऊँची दरों में।
\end{solution}

\begin{solution}{exo:b3:nervous-systems:8}
जो दो न्यूरॉन बिना थके एक-दूसरे का निरोध करें वे कोई स्विच बनाते हैं:
जो पहले दगे वह दूसरे को तब तक दबाए रखता है जब तक चालन चलता रहे, और
बागडोर कभी बदलती ही नहीं। दोलन के लिए कोई ऐसी धीमी प्रक्रिया चाहिए जो
विजेता की पकड़ ढीली करे --- उसके दगने का अनुकूलन, उसकी निरोधी संधि का
अवसादन, या छूटने के बाद हारे हुए की कोई पलटवार-सक्रियता --- और उसी का
काल-अचर आवर्तकाल तय करता है।
\end{solution}

\begin{solution}{exo:b3:nervous-systems:9}
L-डोपा को बड़े उदासीन ऐमीनो अम्लों का वाहक उस अवरोध के पार ले जाता है;
डोपामीन, जो ध्रुवी और आवेशित है, नहीं जाता। रोगी L-डोपा लेते हैं
(परिधीय एंज़ाइम के किसी निरोधक के साथ, ताकि वह पार जाने से पहले बदल न
जाए), और मस्तिष्क की अपनी एंज़ाइम उससे वहीं डोपामीन बना लेती है जहाँ
उसकी ज़रूरत है। \omterm{def:b3:adaptive-immunity:antibody}{प्रतिरक्षी}, \qty{150}{kDa} के, न दृढ़ जोड़ों को पार
करते हैं न किसी वाहक को; पहुँचाने के लिए कोई ऐसी गाड़ी चाहिए जो
पारकोशिकीय परिवहन के लिए किसी ग्राही का अपहरण करे, या
\omterm{prop:b3:nervous-systems:barrier-budget}{प्रमस्तिष्क-मेरु द्रव} में अंतःक्षेपण।
\end{solution}

\begin{solution}{exo:b3:nervous-systems:10}
प्रति एकांक लंबाई झिल्ली-धारा प्रतिरोधी $V/r_{m}$ जमा धारिती
$c_{m}\,\partial V/\partial t$ है; आवेश-संरक्षण देता है
$\partial I/\partial x = -(V/r_{m} + c_{m}\,\partial V/\partial t)$, और
$\partial V/\partial x = -r_{i}I$ से, $\frac{1}{r_{i}}\,\partial^{2}V/
\partial x^{2} = V/r_{m} + c_{m}\,\partial V/\partial t$;
$r_{m}$ से गुणा करने पर: $\lambda^{2}\,\partial^{2}V/\partial x^{2} = V + \tau\,\partial
V/\partial t$। $\partial V/\partial t = 0$ के साथ प्रमेय का समीकरण
लौट आता है। $\lambda$ कोई लंबाई है और $\tau$ कोई काल, इसलिए
$\lambda/\tau$ कोई चाल है --- यानी किसी विक्षोभ के फैलने की गति का
स्वाभाविक मापक्रम।
\end{solution}

\begin{solution}{exo:b3:nervous-systems:11}
$v \propto \sqrt{d}$: चार गुनी चाल के लिए सोलह गुना व्यास चाहिए, यानी
\qty{8}{mm}। क्षेत्रफल $\pi(4\,\text{mm})^{2} \approx \qty{50}{mm^2}$,
$\pi\,(10\,\unit{\micro m})^{2} = \qty{314}{\micro m^2}$ के सामने: यानी
$1.6\times 10^{5}$ का अनुपात। तेज़ अमायलिनी तंतुओं की कोई तंत्रिका
असंभव है --- ऐसा एक ही तंतु पूरी तंत्रिका जितना मोटा होता --- इसलिए
जिस प्राणी को अनेक तेज़ चैनल चाहिए उसे \omterm{def:b3:nervous-systems:cells}{मायलिन} चाहिए, और इसीलिए
कशेरुकियों ने (और कुछ अकशेरुकियों ने स्वतंत्र रूप से) उसे विकसित किया।
\end{solution}

\begin{solution}{exo:b3:nervous-systems:12}
$a = 0.4/30^{0.75} = 0.4/12.8 = 0.031$; \qty{70000}{g} के लिए: $0.031\times
\num{70000}^{0.75} = 0.031\times 4300 \approx \qty{134}{g}$। मानव
\qty{1350}{g} उस पूर्वानुमान का दस गुना है। मस्तिष्क-द्रव्यमान
देह-द्रव्यमान के पीछे इसलिए चलता है कि बड़े शरीर को अधिक संवेदी तथा
प्रेरक न्यूरॉन चाहिए, और इसलिए भी कि न्यूरॉन स्वयं मस्तिष्क के आकार के
साथ बढ़ते हैं; संज्ञान का पीछा प्रांतस्था के न्यूरॉनों की संख्या करती
है, जो इस पर निर्भर है कि वे कितनी सघनता से भरे हैं --- नरवानर दूसरे
स्तनियों से प्रति ग्राम अधिक भरते हैं, और मनुष्य सबसे अधिक।
\end{solution}

\begin{solution}{pb:b3:nervous-systems:1}
\textbf{1.} $\lambda = \sqrt{2\times 10^{4}\,d/400}$: $d =
\qty{1e-4}{cm}$ के लिए, $\sqrt{5\times 10^{-3}} = \qty{0.071}{cm} =
\qty{0.71}{mm}$; \qty{4}{\micro m} के लिए, \qty{1.41}{mm}। $\tau = 2\times
10^{4}\times 10^{-6} = \qty{20}{ms}$।
\textbf{2.} $5\,e^{-0.3/0.71} = \qty{3.3}{mV}$; $5\,e^{-0.3/1.41} =
\qty{4.0}{mV}$।
\textbf{3.} दूरस्थ निवेश क्षीण होकर पहुँचते हैं, इसलिए कोई संधि कहाँ
उतरती है यही उसका भार तय करता है; कोई विभव लगभग $\tau$ में क्षीण हो
जाता है, इसलिए दो निवेश तभी जुड़ते हैं जब वे एक-दूसरे के कोई
\qty{20}{ms} के भीतर पड़ें --- यानी न्यूरॉन \qty{20}{ms} की खिड़की
वाला कोई संपात-संसूचक है।
\textbf{4.} मायलिनी \qty{60}{m/s}; अमायलिनी $1.1\sqrt{10} =
\qty{3.5}{m/s}$; अमायलिनी रहते हुए \qty{60}{m/s} तक पहुँचने के लिए
$d = (60/1.1)^{2}
\approx \qty{3000}{\micro m}$: यानी तीन मिलीमीटर।
\textbf{5.} $\pi(5)^{2} = \qty{79}{\micro m^2}$, $\pi(1500)^{2} =
7\times 10^{6}\,\unit{\micro m^2}$ के सामने: एक की जगह में लगभग
\num{90000} मायलिनी तंत्रिकाक्ष।
\textbf{6.} \qty{60}{m/s} पर \qty{2}{cm} \qty{0.3}{ms} है;
\qty{3.5}{m/s} पर, \qty{5.7}{ms}: यानी कोई \qty{5}{ms} अतिरिक्त विलंब।
इससे भी बुरा, उस नंगी झिल्ली की धारिता बड़ी है और वाहिकाएँ थोड़ी, इसलिए
आख़िरी अक्षत गाँठ से आती धारा उसे देहली तक लाने में चूक सकती है: यानी
चालन-अवरोध।
\textbf{7.} संवेदी $1.2/60 = \qty{20}{ms}$; दो संधियाँ \qty{1}{ms};
प्रेरक $1.1/90 = \qty{12}{ms}$: लगभग \qty{33}{ms} (और तंत्रिका-पेशी
संगम पर मिलीसेकंड भर)।
\textbf{8.} मस्तिष्क तक: $0.6/60 = \qty{10}{ms}$ जमा दो संधियाँ, यानी
\omterm{def:b3:nervous-systems:plan}{मेरुरज्जु} में संकेत के घुसने के \qty{20}{ms} के बाद \qty{11}{ms}: संकेत
प्रांतस्था तक लगभग उसी समय पहुँचता है जब पैर हिलता है; पीड़ा के सचेत
अनुभव के लिए कुछ सौ मिलीसेकंड और संसाधन चाहिए, इसलिए पीड़ा महसूस होने
से पहले ही पैर खींच लिया जाता है।
\textbf{9.} $1.1\sqrt{1} = \qty{1.1}{m/s}$: $1.8/1.1 \approx
\qty{1.6}{s}$। तीखी पहली पीड़ा पतले मायलिनी तंतुओं से दसियों
मिलीसेकंडों में आती है; मंद, जलती दूसरी पीड़ा अमायलिनी C तंतुओं से
सेकंड भर या उससे अधिक बाद।
\textbf{10.} संवेदी $3.2/60 = \qty{53}{ms}$, प्रेरक $3.1/90 =
\qty{34}{ms}$, संधियाँ \qty{1}{ms}: लगभग \qty{90}{ms}। लंबे प्राणियों
के प्रतिवर्त धीमे होते हैं; वे मोटे तंत्रिकाक्षों से और स्थानीय परिपथों
को अधिक सौंपकर इसकी भरपाई करते हैं।
\textbf{11.} दोनों ओर $0.8/80 = \qty{10}{ms}$, यानी \qty{20}{ms}; संधि
\qty{0.5}{ms}, संगम \qty{1}{ms}, बल \qty{5}{ms}: \qty{26.5}{ms}, और
बाकी \qty{30}{ms} तर्कु के अपने पारक्रमण तथा पेशी के भीतर चालन से पूरा
होता है।
\textbf{12.} अमायलिनी तंत्रिकाक्ष सेकंड भर में एक-दो मीटर चालन करते
हैं, इसलिए जो पाश किसी वयस्क में \qty{30}{ms} लेते हैं वे सैकड़ों लेते
हैं, और समय-निर्धारण खराब रहता है; मायलिनन चाल दस से पचास गुना बढ़ा
देता है और समय-निर्धारण सटीक कर देता है, और चलने, पकड़ने तथा बोलने को
यही चाहिए।
\textbf{13.} $20/50\,000 = 4\times 10^{-4}$ मोल/s $= 2.4\times 10^{20}$
ATP प्रति सेकंड; प्रति न्यूरॉन प्रति सेकंड $2.8\times 10^{9}$।
\textbf{14.} $4\times 10^{8} + 7000\times 2\times 10^{4} = 4\times 10^{8}
+ 1.4\times 10^{8} = 5.4\times 10^{8}$ ATP।
\textbf{15.} सब कुछ आवेगों पर: $2.8\times 10^{9}/5.4\times 10^{8} \approx
\qty{5}{Hz}$; आधा: लगभग \qty{2.6}{Hz}।
\textbf{16.} \qty{1}{Hz} पर, $5.4\times 10^{8}/2.8\times 10^{9} \approx
\qty{20}{\%}$ बजट का: कूट विरल है --- अधिकांश न्यूरॉन अधिकांश समय चुप,
और सूचना इसमें कि थोड़े-से कौन दगते हैं और कब।
\textbf{17.} $8.6\times 10^{10}\times 50\times 5.4\times 10^{8} = 2.3
\times 10^{21}$ ATP/s $= 3.9\times 10^{-3}$ मोल/s $\approx \qty{190}{W}$
--- यानी पूरे विश्रामी शरीर से दुगुना।
\textbf{18.} $5\times 16 = \qty{80}{kJ/h} = \qty{22}{W}$: बस इतना ही,
बिना किसी भंडार के। पूर्ति रुकते ही पंप सेकंडों में रुक जाते हैं,
प्रवणताएँ चुक जाती हैं, और लगभग दस सेकंड में चेतना चली जाती है।
\textbf{19.} $a = 0.4/12.8 = 0.031$; \qty{70}{kg}: $0.031\times 4300 =
\qty{134}{g}$; \qty{5000}{kg}: $0.031\times 1.06\times 10^{5} \approx
\qty{3300}{g}$।
\textbf{20.} मनुष्य $1350/134 \approx 10$; हाथी $4800/3300 \approx
1.5$।
\textbf{21.} \omterm{def:b3:nervous-systems:plan}{अनुमस्तिष्क} में --- उनका \qty{98}{\%} --- जो धड़ की चालीस
हज़ार पेशियों का समन्वय करता है। संज्ञान का पीछा प्रांतस्थीय
न्यूरॉन-संख्या करती है, कुल नहीं; \omterm{def:b3:nervous-systems:plan}{अनुमस्तिष्क} किसी विशाल शरीर को चलाने
के प्रेरक काम का पीछा करता है।
\textbf{22.} तंत्रिकाक्ष और द्रुमिकाएँ मस्तिष्क के साथ लंबी होती हैं,
इसलिए हर न्यूरॉन अधिक आयतन घेरता है और न्यूरॉनों की संख्या द्रव्यमान से
धीमे बढ़ती है। नरवानरों में मस्तिष्क बढ़ने पर भी न्यूरॉन का आकार लगभग
अचर रहता है, इसलिए न्यूरॉन-संख्या लगभग द्रव्यमान के अनुपात में बढ़ती
है, और किसी दिए भार का नरवानर मस्तिष्क उसी भार के कृंतक मस्तिष्क से कई
गुना न्यूरॉन रखता है।
\textbf{23.} कोई भुजा स्थानीय परिपथों से अपने ही चूषक तथा जोड़ सँभाल
सकती है, पहुँच सकती है, पकड़ सकती है, भोजन अपने ऊपर आगे सरका सकती है
और अपनी ही त्वचा पकड़ने से बच सकती है; वह देख नहीं सकती, कोई लक्ष्य
चुन नहीं सकती, कोई विभेदन सीख नहीं सकती। जाँच: भुजा की तंत्रिका-रज्जु
मस्तिष्क से काट दीजिए और उस भुजा को भोजन से छुइए --- वह उसे पकड़कर
उधर सरकाती है जहाँ मुँह होता; कटी हुई भुजा भी घंटे भर तक यही करती है।
\textbf{24.} प्रति न्यूरॉन लगभग $23$ संधियाँ, $7000$ के सामने: यानी तीन
सौ गुना कम। वह पूरा मानचित्र बताता है कि कौन न्यूरॉन किसे प्रभावित कर
सकता है, पर हर संयोजन का बल या चिह्न नहीं, न यह कि कौन-से तंत्रिका-नियंत्रक
उन्हें बदलते हैं, न गतिकी --- इसलिए उस कृमि के लिए भी वह तारबंदी-आरेख
व्यवहार का पूर्वानुमान केवल आंशिक रूप से करता है।
\textbf{25.} \qty{1}{\micro m} की द्रुमिका के लिए $\lambda = \qty{0.71}{mm}$;
प्रत्याहरण-विलंब लगभग \qty{33}{ms}; और आधा बजट आवेगों पर लगाने के साथ
औसत दगने-दर लगभग \qty{2.6}{Hz}।
\end{solution}
